\section{Pembangunan Model STARMA}
Model STARMA digunakan untuk menganalisis data yang memiliki ketergantungan spasial dan temporal secara simulta. Identifikasi orde model melibatkan empat parameter utama: $p$ dan $q$ (lag temporal AR dan MA) serta $\lambda$ dan $\mu$ (lag spasial AR dan MA). Orde ini ditentukan menggunakan fungsi STACF \ref{st_ma} dan STPACF \ref{st_ar}, yang merupakan perluasan dari ACF dan PACF dalam konteks spasial-temporal. Setelah orde ditentukan, parameter model diestimasi menggunakan \textbf{Maximum Likelihood Estimation (MLE)}. Diagnostik dilakukan melalui analisis residu untuk memastikan tidak ada autokorelasi spasial-temporal yang tersisa, serta evaluasi akurasi dilakukan dengan metrik seperti RMSE